\documentclass[11pt]{article} 

\usepackage{amsmath}

\usepackage[super]{nth}

\usepackage{epsfig}

\usepackage{graphics}

\usepackage{wasysym}

\usepackage{times}

\usepackage{natbib}

%\usepackage{amssymb}

\setlength{\topmargin}{0in}
\setlength{\textwidth}{6.25in}

\setlength{\oddsidemargin}{.125in}

% 28 lectures total

\usepackage{amscd}

\begin{document}
\title{Quiz \hfill Name : \underline{\hspace{3cm}}}
\date{November 5}
\maketitle
\begin{enumerate}
\item A test for a disease called db is correct 95\% of the time.  That is, if someone has db, the test reports that they do 95\% of the time, if someone does not have db, the test reports that they do not have db 95\% of the time.  The disease of db is rare, affecting 1 in ten thousand people.  
\item[] What is the probability that an individual has db after a positive test (that is, a test that reports db is present)?
\bigskip
\bigskip
\bigskip
\bigskip

\bigskip
\bigskip
\bigskip
\bigskip

\item True or False?
\begin{enumerate}
\item Events $A$ and $A^c$ where $0<P(A)<1$ are never independent.
\item Two independent events are always also conditionally independent.
\item The Total Probability Theorem states that if $A_1,\ldots A_n$ are independent, $P(B)=P(A_1\cap B)+P(A_2\cap B) \ldots+P(A_n\cap B)$/  
\item Events $A$ and $A^c$ partition the sample space.
\item If events Events $A$ and $B^c$ are independent, so are Events $A^c$ and $B^c$
\end{enumerate}


\end{enumerate}
\end{document}